    \chapter*{Notation and Conventions}
    \markboth{Notation and Conventions}{Notation and Conventions}
    \addcontentsline{toc}{chapter}{Notation and Conventions}

    The notation specified below is used throughout this book. When several
    typographic variants of the same letter exist, each variant is assigned a
    distinct role and the variants are not interchanged.

    \section*{General Conventions}

    Scalars are written in ordinary italic letters, such as $a,b,c,\lambda,t$.
    Vectors are written in bold lowercase letters, such as
    $\mathbf{u},\mathbf{v},\mathbf{x}$. Matrices are written in bold uppercase
    letters, such as $\mathbf{A},\mathbf{B},\mathbf{P}$. In particular, the
    identity matrix, a diagonal matrix, and a projection matrix are written as
    $\mathbf{I}_n$, $\mathbf{D}$, and $\mathbf{P}$, respectively. Unless
    otherwise stated, vectors are column vectors.

    Round parentheses are used for ordered tuples and matrix delimiters. Square
    brackets are used for closed intervals and for area notation such as
    $[ABC]$. Braces are used for sets. In set-builder notation, the condition is
    separated by $\mid$, as in $\{x\in A\mid P(x)\}$. Determinants may be
    enclosed by vertical bars because a determinant is a scalar.

    In symbolic statements, logical conjunction, disjunction, and negation are
    written as $\land$, $\lor$, and $\neg$. The words ``and'' and ``or'' are
    used in ordinary prose.

    \vspace{0.5em}

    {\small
    \renewcommand{\arraystretch}{1.18}
    \begin{longtable}{@{}L{0.24\textwidth}L{0.68\textwidth}@{}}
        \toprule
        Notation & Meaning \\
        \midrule
        \endfirsthead

        \toprule
        Notation & Meaning \\
        \midrule
        \endhead

        \bottomrule
        \endlastfoot

        $\N$ & the set of positive integers, $\{1,2,3,\dots\}$ \\
        $\Nzero$ & the set of nonnegative integers, $\{0,1,2,3,\dots\}$ \\
        $\Z$ & the set of integers \\
        $\Q$ & the set of rational numbers \\
        $\R$ & the set of real numbers \\
        $\C$ & the set of complex numbers \\
        $a\coloneqq b$ & $a$ is defined to be $b$ \\
        $\varepsilon$ & a positive quantity in analysis, usually arbitrary \\
        $\epsilon_{ijk}$ & the Levi--Civita symbol; the glyph $\epsilon$ is reserved for this use \\
        $O(g(x))$ & Big-O notation \\
        $o(g(x))$ & little-o notation \\
        $\square$ & the end of a proof or solution \\
    \end{longtable}
    }

    \section*{Theorem-like Environments}

    {\small
    \renewcommand{\arraystretch}{1.18}
    \begin{longtable}{@{}L{0.24\textwidth}L{0.68\textwidth}@{}}
        \toprule
        Label & Usage \\
        \midrule
        \endfirsthead

        \toprule
        Label & Usage \\
        \midrule
        \endhead

        \bottomrule
        \endlastfoot

        Definition & a precise meaning of a term or object \\
        Theorem & a major mathematical result \\
        Proposition & a useful result of more limited scope \\
        Lemma & a supporting result or technical tool \\
        Corollary & a result that follows directly from another result \\
        Formula & a computational identity or closed form \\
        Algorithm & a step-by-step computational procedure \\
        Fact & a verified mathematical or historical fact \\
        Conjecture & a statement believed to be true but not proved in full generality \\
        Hypothesis & a named assumption or proposed statement \\
        Open Problem & a problem whose stated form remains unsolved \\
        Idea & a concise account of a decisive mathematical idea \\
        Strategy & a broadly useful problem-solving method \\
        Tactic & a local, immediately applicable device \\
        Remark & supplementary context or a caution \\
    \end{longtable}
    }

    \section*{Set Theory and Logic}

    Uppercase letters such as $A,B,X,Y$ usually denote sets, while lowercase
    letters such as $a,b,x,y$ usually denote their elements unless the context
    indicates another meaning. The arrow $f:X\to Y$ specifies the domain and
    codomain of a function, while $x\mapsto f(x)$ specifies its action on an
    element.

    \vspace{0.5em}

    {\small
    \renewcommand{\arraystretch}{1.18}
    \begin{longtable}{@{}L{0.27\textwidth}L{0.65\textwidth}@{}}
        \toprule
        Notation & Meaning \\
        \midrule
        \endfirsthead

        \toprule
        Notation & Meaning \\
        \midrule
        \endhead

        \bottomrule
        \endlastfoot

        $x\in A$ & $x$ is an element of $A$ \\
        $x\notin A$ & $x$ is not an element of $A$ \\
        $A\subseteq B$ & $A$ is a subset of $B$ \\
        $A\subsetneq B$ & $A$ is a proper subset of $B$ \\
        $A\supseteq B$ & $A$ is a superset of $B$ \\
        $A\supsetneq B$ & $A$ is a proper superset of $B$ \\
        $A\cup B$ & the union of $A$ and $B$ \\
        $A\cap B$ & the intersection of $A$ and $B$ \\
        $A\setminus B$ & the set difference of $A$ and $B$ \\
        $A^{\mathsf c}$ & the complement of $A$ in the understood ambient set \\
        $\varnothing$ & the empty set \\
        $\Pow(A)$ & the power set of $A$ \\
        $A\times B$ & the Cartesian product of $A$ and $B$ \\
        $A\sqcup B$ & the disjoint union of $A$ and $B$ \\
        $|A|$ & the cardinality of a finite set $A$ \\
        $\{x\in A\mid P(x)\}$ & the set of all $x\in A$ for which $P(x)$ holds \\
        $x\sim y$ & $x$ and $y$ are equivalent under an understood equivalence relation \\
        $A/\sim$ & the set of equivalence classes of $A$ under $\sim$ \\
        $\forall x\in A$ & for every $x\in A$ \\
        $\exists x\in A$ & there exists $x\in A$ \\
        $\exists!x\in A$ & there exists a unique $x\in A$ \\
        $\nexists x\in A$ & there does not exist $x\in A$ \\
        $P\Rightarrow Q$ & $P$ implies $Q$ \\
        $P\Leftarrow Q$ & $Q$ implies $P$ \\
        $P\Leftrightarrow Q$ & $P$ if and only if $Q$ \\
        $P\land Q$ & logical conjunction \\
        $P\lor Q$ & logical disjunction \\
        $\neg P$ & logical negation \\
        $f:X\to Y$ & a function from $X$ to $Y$ \\
        $x\mapsto f(x)$ & the elementwise rule of a function \\
        $\dom f$ & the domain of $f$ \\
        $\codom f$ & the codomain of $f$ \\
        $\id_X$ & the identity map on $X$ \\
        $g\circ f$ & composition of functions \\
        $f|_A$ & restriction of $f$ to $A$ \\
        $f:X\hookrightarrow Y$ & an injective map \\
        $f:X\twoheadrightarrow Y$ & a surjective map \\
        $f:X\xrightarrow{\sim}Y$ & an isomorphism \\
        $X\cong Y$ & $X$ and $Y$ are isomorphic; in topology, this means homeomorphic \\
        $f(A)$ & the image of $A$ under $f$ \\
        $f^{-1}(B)$ & the inverse image, or preimage, of $B$ \\
        $\indicatorbb{A}(x)$ & the indicator function of $A$ \\
        $\aleph_0$ & the cardinality of a countably infinite set \\
        $\mathfrak{c}$ & the cardinality of the continuum \\
    \end{longtable}
    }

    The symbol $\indicatorbb{A}$ is blackboard bold and does not denote a
    vector. The all-ones vector is $\ones$. Boldface is otherwise reserved for
    vectors and matrices.

    \section*{Calculus and Analysis}

    Logarithms to bases $e$, $10$, and $2$ are written as $\ln x$, $\lg x$,
    and $\lb x$, respectively.

    The coordinate conventions are
    \[
        \begin{aligned}
            &\text{Cartesian:} &&(x,y),\quad(x,y,z),\\
            &\text{polar:} &&(r,\theta),\\
            &\text{cylindrical:} &&(r,\theta,z),\\
            &\text{spherical:} &&(\rho,\theta,\phi).
        \end{aligned}
    \]
    In polar, cylindrical, and spherical coordinates, $\theta$ is the azimuthal
    angle measured counterclockwise from the positive $x$-axis in the
    $xy$-plane. In spherical coordinates, $\phi$ is the polar angle measured
    from the positive $z$-axis. The glyph $\phi$ is reserved for this spherical
    polar angle. The glyph $\varphi$ is used for functions, homomorphisms,
    Euler's totient function, and the golden ratio whenever one of these objects
    is introduced.

    \vspace{0.5em}

    {\small
    \renewcommand{\arraystretch}{1.18}
    \begin{longtable}{@{}L{0.27\textwidth}L{0.65\textwidth}@{}}
        \toprule
        Notation & Meaning \\
        \midrule
        \endfirsthead

        \toprule
        Notation & Meaning \\
        \midrule
        \endhead

        \bottomrule
        \endlastfoot

        $f'(a)$ & the first derivative of $f$ at $a$ \\
        $f^{(n)}(a)$ & the $n$th derivative of $f$ at $a$ \\
        $\displaystyle\frac{\dd y}{\dd x}$ & the derivative of $y$ with respect to $x$, with upright $\dd$ \\
        $\displaystyle\int_a^b f(x)\dd x$ & the definite integral of $f$ over $[a,b]$ \\
        $\ln x$ & the logarithm to base $e$ \\
        $\lg x$ & the logarithm to base $10$ \\
        $\lb x$ & the logarithm to base $2$ \\
        $\log_a x$ & the logarithm to the explicitly specified base $a$ \\
        $\arcsin x$ & the principal inverse sine, with values in $[-\pi/2,\pi/2]$ \\
        $\arccos x$ & the principal inverse cosine, with values in $[0,\pi]$ \\
        $\arctan x$ & the principal inverse tangent, with values in $(-\pi/2,\pi/2)$ \\
        $C^k(I)$ & the space of $k$ times continuously differentiable functions on $I$ \\
        $C^\infty(I)$ & the space of smooth functions on $I$ \\
        $\displaystyle\frac{\partial f}{\partial x_i}$ & the partial derivative of $f$ with respect to $x_i$ \\
        $(x,y)$, $(x,y,z)$ & Cartesian coordinates \\
        $(r,\theta)$ & polar coordinates \\
        $(r,\theta,z)$ & cylindrical coordinates \\
        $(\rho,\theta,\phi)$ & spherical coordinates \\
        $\nabla f$ & the gradient of $f$ \\
        $\nabla\cdot\mathbf{F}$ & the divergence of $\mathbf{F}$ \\
        $\nabla\times\mathbf{F}$ & the curl of $\mathbf{F}$ \\
        $\Delta f$ & the Laplacian of $f$ \\
        $D\mathbf{F}(\mathbf{a})$ & the total derivative, or Jacobian matrix, of $\mathbf{F}$ at $\mathbf{a}$ \\
        $\Jac(\mathbf{F})$ & the Jacobian determinant of $\mathbf{F}$ \\
        $\Hess f$ & the Hessian matrix of $f$ \\
        $\mathbf{r}(t)$ & a parametrized curve \\
        $\dd s$ & an element of arc length \\
        $\displaystyle\int_C\mathbf{F}\cdot\dd\mathbf{r}$ & a line integral along $C$ \\
        $\displaystyle\iint_S\mathbf{F}\cdot\mathbf{n}\dd S$ & a flux integral across $S$ \\
        $B_r(x)$ & the open ball of radius $r$ centered at $x$ \\
        $\overline{B}_r(x)$ & the closed ball of radius $r$ centered at $x$ \\
        $x_n\to x$ & convergence of a sequence \\
        $a_n\nearrow L$ & monotone increasing convergence to $L$ \\
        $a_n\searrow L$ & monotone decreasing convergence to $L$ \\
        $\limsup a_n$ & the limit superior of $(a_n)$ \\
        $\liminf a_n$ & the limit inferior of $(a_n)$ \\
        $f_n\ptwto f$ & pointwise convergence \\
        $f_n\unifto f$ & uniform convergence \\
        $f_n\aeto f$ & convergence almost everywhere \\
        $\|f\|_\infty$ & the supremum norm of $f$ \\
        $z=x+iy$ & a complex number with real part $x$ and imaginary part $y$ \\
        $\Real z$ & the real part of $z$ \\
        $\Imag z$ & the imaginary part of $z$ \\
        $\overline{z}$ & the complex conjugate of $z$ \\
        $|z|$ & the modulus of $z$ \\
        $\arg z$ & an argument of $z$ \\
        $\displaystyle\oint_C f(z)\dd z$ & a contour integral around a closed curve $C$ \\
        $\Res(f;a)$ & the residue of $f$ at $a$ \\
    \end{longtable}
    }

    \section*{Linear Algebra}

    Vectors and matrices are bold; scalars and linear maps are ordinary italic.

    \vspace{0.5em}

    {\small
    \renewcommand{\arraystretch}{1.18}
    \begin{longtable}{@{}L{0.27\textwidth}L{0.65\textwidth}@{}}
        \toprule
        Notation & Meaning \\
        \midrule
        \endfirsthead

        \toprule
        Notation & Meaning \\
        \midrule
        \endhead

        \bottomrule
        \endlastfoot

        $\mathbf{A}=(a_{ij})$ & the matrix whose $(i,j)$-entry is $a_{ij}$ \\
        $\mathbf{A}^{\mathsf T}$ & the transpose of $\mathbf{A}$ \\
        $\mathbf{A}^{\mathsf H}$ & the conjugate transpose of $\mathbf{A}$ \\
        $\mathbf{A}^{-1}$ & the inverse of $\mathbf{A}$ \\
        $\mathbf{I}_n$ & the $n\times n$ identity matrix \\
        $\mathbf{0}$ & a zero vector or zero matrix, with dimension clear from context \\
        $\mathbf{D}$ & a diagonal matrix \\
        $\ones$ & the all-ones vector \\
        $\mathbf{J}_{m,n}$, $\mathbf{J}_n$ & the $m\times n$ and $n\times n$ all-ones matrices \\
        $\mathbf{P}_\sigma$ & the permutation matrix associated with $\sigma$ \\
        $\det(\mathbf{A})$ & the determinant of $\mathbf{A}$ \\
        $\rk(\mathbf{A})$ & the rank of $\mathbf{A}$ \\
        $\tr(\mathbf{A})$ & the trace of $\mathbf{A}$ \\
        $p_{\mathbf{A}}(\lambda)$ & the characteristic polynomial of $\mathbf{A}$ \\
        $m_{\mathbf{A}}(\lambda)$ & the minimal polynomial of $\mathbf{A}$ \\
        $\rho(\mathbf{A})$ & the spectral radius of $\mathbf{A}$ \\
        $\sigma_i(\mathbf{A})$ & the $i$th singular value of $\mathbf{A}$ \\
        $M_{ij}$ & the minor of the entry $a_{ij}$ \\
        $C_{ij}=(-1)^{i+j}M_{ij}$ & the cofactor of the entry $a_{ij}$ \\
        $\adj(\mathbf{A})$ & the adjugate of $\mathbf{A}$ \\
        $\Col(\mathbf{A})$ & the column space of $\mathbf{A}$ \\
        $\Null(\mathbf{A})$ & the null space of $\mathbf{A}$ \\
        $\Row(\mathbf{A})$ & the row space of $\mathbf{A}$ \\
        $\ker T$ & the kernel of a linear map $T$ \\
        $\image T$ & the image of a linear map $T$ \\
        $\mathbf{e}_i$ & the $i$th standard basis vector \\
        $\Span(\mathbf{v}_1,\dots,\mathbf{v}_k)$ & the span of the listed vectors \\
        $U\oplus V$ & the direct sum of subspaces $U$ and $V$ \\
        $\langle\mathbf{u},\mathbf{v}\rangle$ & the inner product of $\mathbf{u}$ and $\mathbf{v}$ \\
        $\|\mathbf{v}\|$ & the norm of $\mathbf{v}$ \\
        $\proj_{\mathbf{v}}\mathbf{x}$ & the orthogonal projection of $\mathbf{x}$ onto $\Span(\mathbf{v})$ \\
        $\mathbf{u}\perp\mathbf{v}$ & $\mathbf{u}$ and $\mathbf{v}$ are orthogonal \\
        $\mathbf{u}\times\mathbf{v}$ & the cross product in $\R^3$ \\
    \end{longtable}
    }

    \section*{Differential Equations}

    {\small
    \renewcommand{\arraystretch}{1.18}
    \begin{longtable}{@{}L{0.27\textwidth}L{0.65\textwidth}@{}}
        \toprule
        Notation & Meaning \\
        \midrule
        \endfirsthead

        \toprule
        Notation & Meaning \\
        \midrule
        \endhead

        \bottomrule
        \endlastfoot

        $y'$ & the first derivative of $y$ with respect to the independent variable \\
        $y''$ & the second derivative of $y$ \\
        $\dot{x}$ & the first derivative of $x$ with respect to time $t$ \\
        $\ddot{x}$ & the second derivative of $x$ with respect to time $t$ \\
        $y_h$ & the homogeneous part of a solution \\
        $y_p$ & a particular solution \\
        $c_1,c_2,\dots$ & arbitrary constants in a general solution \\
        $W(y_1,y_2)(t)$ & the Wronskian of $y_1$ and $y_2$ \\
        $\mathcal{L}\{f(t)\}(s)$ & the Laplace transform of $f(t)$ \\
        $\widehat{f}(\omega)$ & the Fourier transform under the convention stated where it is used \\
    \end{longtable}
    }

    \section*{Combinatorics and Discrete Mathematics}

    The symbol $\multichoose{n}{k}$ is reserved for combinations with
    repetition.

    \vspace{0.5em}

    {\small
    \renewcommand{\arraystretch}{1.18}
    \begin{longtable}{@{}L{0.27\textwidth}L{0.65\textwidth}@{}}
        \toprule
        Notation & Meaning \\
        \midrule
        \endfirsthead

        \toprule
        Notation & Meaning \\
        \midrule
        \endhead

        \bottomrule
        \endlastfoot

        $[n]$ & the set $\{1,2,\dots,n\}$ \\
        $n!$ & factorial \\
        $\binom{n}{k}$ & the number of $k$-element subsets of an $n$-element set \\
        $\multichoose{n}{k}$ & the number of size-$k$ multisets chosen from an $n$-element set \\
        $P(n,k)$ & the number of $k$-permutations of an $n$-element set \\
        $\stirlingsecond{n}{k}$ & the Stirling number of the second kind \\
        $a_n$ & the $n$th term of a sequence \\
        $\Delta a_n$ & the forward difference $a_{n+1}-a_n$ \\
        $S_n=\displaystyle\sum_{k=1}^{n}a_k$ & the $n$th partial sum \\
        $A(x)=\displaystyle\sum_{n\geq0}a_nx^n$ & an ordinary generating function \\
        $[x^n]A(x)$ & the coefficient of $x^n$ in $A(x)$ \\
        $\omega_n=e^{2\pi i/n}$ & a primitive $n$th root of unity \\
        $\sigma=(a_1\ a_2\ \cdots\ a_k)$ & cycle notation for a permutation \\
        $\inv(\sigma)$ & the inversion number of $\sigma$ \\
        $\Fix(g)$ & the set of objects fixed by $g$ \\
        $G=(V,E)$ & a graph with vertex set $V$ and edge set $E$ \\
        $\deg(v)$ & the degree of the vertex $v$ \\
        $K_n$ & the complete graph on $n$ vertices \\
        $C_n$ & the cycle graph on $n$ vertices \\
        $P_n$ & the path graph on $n$ vertices \\
    \end{longtable}
    }

    \section*{Probability and Statistics}

    Random variables are usually denoted by uppercase letters such as $X,Y,Z$;
    events are denoted by uppercase letters such as $A,B,E$.

    \vspace{0.5em}

    {\small
    \renewcommand{\arraystretch}{1.18}
    \begin{longtable}{@{}L{0.27\textwidth}L{0.65\textwidth}@{}}
        \toprule
        Notation & Meaning \\
        \midrule
        \endfirsthead

        \toprule
        Notation & Meaning \\
        \midrule
        \endhead

        \bottomrule
        \endlastfoot

        $\Prob(A)$ & the probability of the event $A$ \\
        $\Prob(A\mid B)$ & the conditional probability of $A$ given $B$ \\
        $\Exp[X]$ & the expectation of $X$ \\
        $\Var(X)$ & the variance of $X$ \\
        $\Cov(X,Y)$ & the covariance of $X$ and $Y$ \\
        $X\sim\Bin(n,p)$ & a binomial distribution \\
        $X\sim\Poi(\lambda)$ & a Poisson distribution \\
        $X\sim\Normal(\mu,\sigma^2)$ & a normal distribution with mean $\mu$ and variance $\sigma^2$ \\
        $X_n\asto X$ & almost-sure convergence \\
        $X_n\probto X$ & convergence in probability \\
        $X_n\Lpto{p}X$ & convergence in $L^p$ \\
        $X_n\distto X$ & convergence in distribution \\
    \end{longtable}
    }

    \section*{Geometry}

    Uppercase letters such as $A,B,C,D$ denote points. When the usual triangle
    notation is adopted, the corresponding lowercase letters $a,b,c$ denote the
    side lengths opposite $A,B,C$, respectively. For a triangle, $R$, $r$, $s$,
    and $[ABC]$ denote the circumradius, inradius, semiperimeter, and area, respectively.

    \vspace{0.5em}

    {\small
    \renewcommand{\arraystretch}{1.18}
    \begin{longtable}{@{}L{0.24\textwidth}L{0.68\textwidth}@{}}
        \toprule
        Notation & Meaning \\
        \midrule
        \endfirsthead

        \toprule
        Notation & Meaning \\
        \midrule
        \endhead

        \bottomrule
        \endlastfoot

        $\triangle ABC$ & the triangle with vertices $A,B,C$ \\
        $\angle ABC$ & the angle with vertex $B$ \\
        $\overline{AB}$ & the segment joining $A$ and $B$ \\
        $AB$ & the length of $\overline{AB}$ \\
        $[ABC]$ & the area of $\triangle ABC$ \\
        $\conv(S)$ & the convex hull of $S$ \\
        $R$ & the circumradius of a triangle \\
        $r$ & the inradius of a triangle \\
        $s$ & the semiperimeter of a triangle \\
        $I$ & the number of interior lattice points in Pick's Theorem \\
        $B$ & the number of boundary lattice points in Pick's Theorem \\
        $\parallel$ & parallel \\
        $\perp$ & perpendicular \\
    \end{longtable}
    }

    \section*{Number Theory}

    Unless otherwise stated, divisibility and congruence statements are made in
    $\Z$.

    \vspace{0.5em}

    {\small
    \renewcommand{\arraystretch}{1.18}
    \begin{longtable}{@{}L{0.27\textwidth}L{0.65\textwidth}@{}}
        \toprule
        Notation & Meaning \\
        \midrule
        \endfirsthead

        \toprule
        Notation & Meaning \\
        \midrule
        \endhead

        \bottomrule
        \endlastfoot

        $a\mid b$ & $a$ divides $b$ \\
        $a\nmid b$ & $a$ does not divide $b$ \\
        $\gcd(a,b)$ & the greatest common divisor of $a$ and $b$ \\
        $\lcm(a,b)$ & the least common multiple of $a$ and $b$ \\
        $a\equiv b\pmod n$ & $a$ is congruent to $b$ modulo $n$ \\
        $[a]_n$ & the residue class of $a$ modulo $n$ \\
        $\Z/n\Z$ & the ring of residue classes modulo $n$ \\
        $(\Z/n\Z)^\times$ & the multiplicative group of units modulo $n$ \\
        $a^{-1}\pmod n$ & the multiplicative inverse of $a$ modulo $n$ \\
        $\left(\frac{a}{p}\right)$ & the Legendre symbol \\
        $\varphi(n)$ & Euler's totient function \\
        $\ord_n(a)$ & the multiplicative order of $a$ modulo $n$ \\
        $\nu_p(n)$ & the exponent of $p$ in the prime factorization of $n$ \\
        $\tau(n)$ & the number of positive divisors of $n$ \\
        $\sigma(n)$ & the sum of the positive divisors of $n$ \\
        $\mu(n)$ & the M\"obius function \\
        $p$ & usually a prime number in number-theoretic contexts \\
        $\mathbb{F}_q$ & the finite field with $q$ elements \\
    \end{longtable}
    }

    \section*{Algebra}

    Capital letters such as $G,H$ denote groups, $R,S$ denote rings, and
    $F,K,L$ denote fields; lowercase letters denote their elements unless the
    context requires otherwise.

    \vspace{0.5em}

    {\small
    \renewcommand{\arraystretch}{1.18}
    \begin{longtable}{@{}L{0.27\textwidth}L{0.65\textwidth}@{}}
        \toprule
        Notation & Meaning \\
        \midrule
        \endfirsthead

        \toprule
        Notation & Meaning \\
        \midrule
        \endhead

        \bottomrule
        \endlastfoot

        $G$ & usually a group \\
        $S_n$ & the symmetric group on $n$ symbols \\
        $A_n$ & the alternating group on $n$ symbols \\
        $D_{2n}$ & the dihedral group of order $2n$ \\
        $e$ & the identity element of a group \\
        $|G|$ & the order of a finite group $G$ \\
        $\langle S\rangle$ & the subgroup generated by $S$ \\
        $H\leq G$ & $H$ is a subgroup of $G$ \\
        $H\trianglelefteq G$ & $H$ is a normal subgroup of $G$ \\
        $G/H$ & the quotient group of $G$ by $H$ \\
        $\ker\varphi$ & the kernel of a homomorphism $\varphi$ \\
        $\image\varphi$ & the image of a homomorphism $\varphi$ \\
        $\Aut(G)$ & the automorphism group of $G$ \\
        $G\cdot x$ & the orbit of $x$ under a group action \\
        $G_x$ & the stabilizer of $x$ under a group action \\
        $Z(G)$ & the center of $G$ \\
        $C_G(x)$ & the centralizer of $x$ in $G$ \\
        $\GL_n(F)$ & the general linear group over $F$ \\
        $\SL_n(F)$ & the special linear group over $F$ \\
        $R$ & usually a ring \\
        $I\trianglelefteq R$ & $I$ is an ideal of $R$ \\
        $R/I$ & the quotient ring of $R$ by $I$ \\
        $F$ & a base field in an algebraic context \\
        $K$ & a field extension of $F$, or a second field when required \\
        $F[x]$ & the polynomial ring over $F$ \\
        $\deg f$ & the degree of a polynomial $f$ \\
        $\Disc(f)$ & the discriminant of a polynomial $f$ \\
        $\characteristic F$ & the characteristic of $F$ \\
        $\Gal(L/K)$ & the Galois group of the extension $L/K$ \\
    \end{longtable}
    }

    \section*{Topology}

    {\small
    \renewcommand{\arraystretch}{1.18}
    \begin{longtable}{@{}L{0.27\textwidth}L{0.65\textwidth}@{}}
        \toprule
        Notation & Meaning \\
        \midrule
        \endfirsthead

        \toprule
        Notation & Meaning \\
        \midrule
        \endhead

        \bottomrule
        \endlastfoot

        $(X,\mathcal{T})$ & a topological space \\
        $\interior A$ & the interior of $A$ \\
        $\overline{A}$ & the closure of $A$ \\
        $\partial A$ & the boundary of $A$ \\
        $\pi_1(X,x_0)$ & the fundamental group of $X$ based at $x_0$ \\
        $\chi(X)$ & the Euler characteristic of $X$ \\
    \end{longtable}
    }

    \section*{Miscellaneous Topics}

    \noindent\textbf{Mathematical Constants}

    \vspace{0.5em}

    The table contains constants that occur frequently across several areas of
    mathematics. A defining expression is used whenever a single-letter symbol
    is not needed elsewhere in the book.

    \vspace{0.5em}

    {\small
    \renewcommand{\arraystretch}{1.18}
    \begin{longtable}{@{}L{0.12\textwidth}L{0.22\textwidth}L{0.38\textwidth}L{0.20\textwidth}@{}}
        \toprule
        Symbol & Name & Meaning or definition & Exact value or approximation \\
        \midrule
        \endfirsthead

        \toprule
        Symbol & Name & Meaning or definition & Exact value or approximation \\
        \midrule
        \endhead

        \bottomrule
        \endlastfoot

        $i$
        & Imaginary unit
        & The imaginary unit, defined by the relation $i^2=-1$.
        & $i^2=-1$ \\

        $\sqrt{2}$
        & Pythagoras constant
        & The positive length of the diagonal of a unit square.
        & $1.41421356237\ldots$ \\

        $e$
        & Euler's number
        & The base of the natural logarithm:
          $e=\sum_{n=0}^{\infty}1/n!$.
        & $2.71828182846\ldots$ \\

        $\pi$
        & Pi
        & The ratio of the circumference of a Euclidean circle to its diameter.
        & $3.14159265359\ldots$ \\

        $\varphi$
        & Golden ratio
        & The positive solution of $x^2=x+1$.
        & $\dfrac{1+\sqrt{5}}{2}\approx1.61803398875\ldots$ \\

        $\gamma$
        & Euler--Mascheroni constant
        & $\displaystyle
          \gamma=\lim_{n\to\infty}
          \left(\sum_{k=1}^{n}\frac{1}{k}-\ln n\right)$.
        & $0.577215664902\ldots$ \\

        $\zeta(3)$
        & Apéry's constant
        & The value of the Riemann zeta function at $3$:
          $\displaystyle \sum_{n=1}^{\infty}1/n^3$.
        & $1.20205690316\ldots$ \\

        $G$
        & Catalan's constant
        & $\displaystyle
          G=\sum_{n=0}^{\infty}\frac{(-1)^n}{(2n+1)^2}$.
        & $0.915965594177\ldots$ \\
    \end{longtable}
    }

    \vspace{1em}

    \noindent\textbf{Other Notation}

    \vspace{0.5em}

    {\small
    \renewcommand{\arraystretch}{1.18}
    \begin{longtable}{@{}L{0.27\textwidth}L{0.65\textwidth}@{}}
        \toprule
        Notation & Meaning \\
        \midrule
        \endfirsthead

        \toprule
        Notation & Meaning \\
        \midrule
        \endhead

        \bottomrule
        \endlastfoot

        $\lfloor x\rfloor$ & the floor of $x$ \\
        $\lceil x\rceil$ & the ceiling of $x$ \\
        $\{x\}$ & the fractional part $x-\lfloor x\rfloor$ \\
        RH & the Riemann Hypothesis, after it has been introduced \\
        P vs NP & the P versus NP problem, after it has been introduced \\
    \end{longtable}
    }
